\documentclass[../Main.tex]{subfiles}
\begin{document}
  \chapter{2022-2023学年微积分（一）（下）期末考试}

  \section{单项选择题(每小题 3 分，共 18 分)}
  \begin{enumerate}
    \item 在空间直角坐标系中，方程 \( z = \sqrt{x^2 + y^2 - 1} \) 表示（\qquad）.

      \begin{tasks}[label=\textbf{\Alph*.}, label-width=1.5em, column-sep=1em](2)
        \task 半球面
        \task 柱面
        \task 锥面
        \task 单叶双曲面
      \end{tasks}

    \item 设 \( z = f(x, y) \) 在点 \( (1,1) \) 处有 \( f_x(1,1) = f_y(1,1) = 1 \)，则必有（\qquad）.

      \begin{tasks}[label=\textbf{\Alph*.}, label-width=1.5em, column-sep=1em](1)
        \task \( \lim\limits_{\substack{x \to 1 \\ y \to 1}} f(x, y) \) 存在
        \task \( \mathrm{d}z|_{(1,1)} = \mathrm{d}x + \mathrm{d}y \)
        \task \( \lim\limits_{x \to 1} f(x, 1) \) 及 \( \lim\limits_{y \to 1} f(1, y) \) 都存在
        \task \( f(x, y) \) 在点 \( (1,1) \) 沿方向 \( \boldsymbol{n} = \{\cos\theta, \sin\theta\} \) 的方向导数存在
      \end{tasks}

    \item 设函数 \( z = f(x, y) \) 的全微分为 \( \mathrm{d}z = 2x \mathrm{d}x + 3y \mathrm{d}y \)，则点 \( (0,0) \)（\qquad）.

      \begin{tasks}[label=\textbf{\Alph*.}, label-width=1.5em, column-sep=1em](2)
        \task 不是 \( f(x, y) \) 的连续点
        \task 不是 \( f(x, y) \) 的驻点
        \task 是 \( f(x, y) \) 的极大值点
        \task 是 \( f(x, y) \) 的极小值点
      \end{tasks}

    \item 设 \( \Omega \) 是由锥面 \( z = \sqrt{x^2 + y^2} \) 和平行 \( z = 1 \) 所围成的空间区域，将 \( I = \displaystyle \iiint \limits_{\Omega} f(\sqrt{x^2 + y^2 + z}) \mathrm{d}x\,\mathrm{d}y\,\mathrm{d}z \) 化为柱坐标系下的累次积分，下列结果正确的是（\qquad）.

      \begin{tasks}[label=\textbf{\Alph*.}, label-width=1.5em, column-sep=1em](1)
        \task \( I = 2\pi \displaystyle\int_0^1 r \mathrm{d}r \displaystyle\int_0^1 f(\sqrt{r^2 + z}) \mathrm{d}z \)
        \task \( I = 2\pi \displaystyle\int_0^1 r \mathrm{d}r \displaystyle \int_r^1 f(\sqrt{r^2 + z}) \mathrm{d}z \)
        \task \( I = 2\pi \displaystyle\int_0^1 \mathrm{d}r \displaystyle\int_r^1 f(\sqrt{r^2 + z}) \mathrm{d}z \)
        \task \( I = 2\pi \displaystyle\int_0^1 \mathrm{d}z \displaystyle \int_0^z f(\sqrt{r^2 + z}) \mathrm{d}r \)
      \end{tasks}

    \item 设 \( S \) 是上半球面 \( x^2 + y^2 + z^2 = R^2 \) (\( z \geq 0, R > 0 \))，\( abc \neq 0 \)，则 \( \displaystyle\iint \limits_S (ax + by + cz) \mathrm{d}S = \)（\qquad）.

      \begin{tasks}[label=\textbf{\Alph*.}, label-width=1.5em, column-sep=1em](2)
        \task \( c\pi R^2 \)
        \task \( \dfrac{1}{4}c\pi R^3 \)
        \task \( c\pi R^3 \)
        \task \( (a + b + c)\pi R^2 \)
      \end{tasks}
    \item 下列命题中，正确的是（\qquad）.

      \begin{tasks}[label=\textbf{\Alph*.}, label-width=1.5em, column-sep=1em](1)
        \task 若 \( a_n \leq b_n \leq c_n \) 且 \( \sum \limits_{n=1}^{\infty} a_n \) 与 \( \sum \limits_{n=1}^{\infty} c_n \) 收敛，则 \( \sum \limits_{n=1}^{\infty} b_n \) 收敛
        \task 若正项级数 \( \sum \limits_{n=1}^{\infty} a_n \) 满足 \( \dfrac{a_{n+1}}{a_n} < 1 \)，则 \( \sum \limits_{n=1}^{\infty} a_n \) 收敛
        \task 若 \( \lim \limits_{n \to \infty} \dfrac{u_n}{v_n} = 1 \)，则 \( \sum \limits_{n=1}^{\infty} u_n \) 与 \( \sum \limits_{n=1}^{\infty} v_n \) 同敛散
        \task 若 \( \sum \limits_{n=1}^{\infty} a_n \) 收敛，则 \( \sum \limits_{n=1}^{\infty} \dfrac{(-1)^n a_n}{\sqrt{n}} \) 收敛
      \end{tasks}
      
  \end{enumerate}

  \section{填空题(每小题 4 分，共 16 分)}
  \begin{enumerate}
    \setcounter{enumi}{6}
    \item 设曲线 \( \Gamma: \begin{cases} x^2 + y^2 + z^2 = 1 \\ x + y + z = 0 \end{cases} \)，则 \( \displaystyle \oint_{\Gamma} (x^2 + 2y^2 + 3z) \mathrm{d}s = \underline{\hspace{5em}} \).
  
    \item \( z = x + 3y + \ln(1 + x^2 + y^2) \) 在点 \( (0,0,0) \) 的切平面方程为 \(\underline{\hspace{5em}}\).
    
    \item 设 \( u = \ln \sqrt{x^2 + y^2 + z^2} \)，则 \( \text{div}(\textbf{grad}u) \Big|_{(1,-1,2)} = \underline{\hspace{5em}} \).
    
    \item 设 \( f(x) = \begin{cases} x, & -\pi \leq x < 0, \\ 1, & x = 0, \\ 2x, & 0 < x < \pi \end{cases} \)，将 \( f(x) \) 展开成以 \( 2\pi \) 为周期的傅里叶级数，则 \( f(x) \) 的傅里叶级数的和函数在 \( x = -3\pi \) 处的值为 \(\underline{\hspace{5em}}\).
    
  \end{enumerate}

  \section{基本计算题(每小题 7 分，共 42 分)}
  \begin{enumerate}
    \setcounter{enumi}{10}
    \item 求过点 \( P(2,1,3) \) 且与 \( L_1: \dfrac{x-1}{2} = \dfrac{y-2}{1} = \dfrac{z-2}{-1} \) 垂直的平面方程，并求该平面与 \( L_1 \) 的交点.
      \vspace{9em}

    \item 设 \( z = e^{-x} \sin \dfrac{y}{x} \)，求 \( \left. \dfrac{\partial z}{\partial y} \right|_{(2,\pi)} \)，\( \left. \dfrac{\partial^2 z}{\partial y \partial x} \right|_{(2,\pi)} \).
      \vspace{9em}

    \item 求曲面 \( z = xy \) 包含在圆柱面 \( x^2 + y^2 = 1 \) 内那部分面积 \( S \).
      \vspace{9em}
    \item 求 \( I = \displaystyle \iiint \limits_{\Omega} z^2 \mathrm{d}x\,\mathrm{d}y\,\mathrm{d}z \)，其中 \( \Omega \) 由 \( \begin{cases} y^2 + \dfrac{z^2}{2} = 1 \\ x = 0 \end{cases} \) 绕 \( z \) 轴旋转一周所形成的曲面所围成的立体.
      \vspace{9em}

    \item 设 \( S \) 是上半球面 \( z = \sqrt{R^2 - x^2 - y^2} \) (\( R > 0 \))，取上侧，求 \( I = \displaystyle \iint \limits_S \dfrac{(x^2 + y) \mathrm{d}y\,\mathrm{d}z + z \mathrm{d}x\,\mathrm{d}y}{\sqrt{x^2 + y^2 + z^2}} \).
      \vspace{9em}

    \item 求幂级数 \( \sum \limits_{n=1}^{\infty} \dfrac{2n + 1}{n!} x^{2n} \) 的收敛域与和函数 \( S(x) \).
      \vspace{9em}
  \end{enumerate}

  \section{综合题(每小题 7 分，共 14 分)}

  \begin{enumerate}
    \setcounter{enumi}{16}
    \item 设 \( L \) 是 \( xOy \) 面上任意的光滑曲线，\( f(x) \) 具有二阶连续的导数，且 \( f(0) = 4 \)，\( f'(0) = 3 \)，若曲线积分 \( \displaystyle\int_L (-xe^{-x} + f''(x))y \mathrm{d}x + f'(x) \mathrm{d}y \) 与路径无关，求 \( f(x) \).
      \vspace{11em}

    \item 求 \( f(x, y) = x^2 - y^2 + 2 \) 在椭圆域 \( D = \{(x, y) \mid 4x^2 + y^2 \leq 4\} \) 上的最大值和最小值.
      \vspace{11em}
  \end{enumerate}

  \section{证明题(每小题 5 分，共 10 分)}

  \begin{enumerate}
    \setcounter{enumi}{18}
    \item 设 \( f(x) \) 在 \([0,1]\) 上连续，且 \( \displaystyle \int_0^1 f(x) dx = A \)，证明：\( \displaystyle \int_0^1 \mathrm{d}x \displaystyle \int_x^1 f(x)f(y) \mathrm{d}y = \dfrac{1}{2}A^2 \).
      \vspace{11em}

    \item 设数列 \( \{a_n\} \) 满足 \( a_1 = 1 \)，\( 2a_{n+1} = a_n + \sqrt{a_n^2 + u_n} \)，其中 \( u_n > 0 \)，\( \{a_n\} \) 有上界，证明 \( \sum \limits_{n=1}^{\infty} u_n \) 收敛.
      \vspace{11em}
  \end{enumerate}
\chapter{2022-2023学年微积分（一）（下）期末考试参考答案}
  \section{单项选择题(每小题 3 分，共 18 分)}
  \begin{enumerate}
    \item \textbf{Solution}. D.

      方程
      \[
        z=\sqrt{x^2+y^2-1}
      \]
      等价于
      \[
        x^2+y^2-z^2=1,
      \]
      表示单叶双曲面的上半部分.

    \item \textbf{Solution}. C.

      在点 \(P(0,1,1)\) 处，
      \[
        F_x= y+ze^{xz}=2,\quad F_y=x-\frac{z}{y}=-1,\quad F_z=-\ln y+xe^{xz}=0.
      \]
      因 \(F_y\neq 0\)，可把 \(y\) 看作 \(x,z\) 的隐函数，且
      \[
        \frac{\partial y}{\partial x}=-\frac{F_x}{F_y}=2.
      \]
      因此 C 中的 \(-2\) 错误.

    \item \textbf{Solution}. D.

      区域 \(D\) 关于原点中心对称，而 \(xy+\cos x\sin y\) 关于 \((x,y)\mapsto(-x,-y)\) 为奇函数，所以
      \[
        \iint_D(xy+\cos x\sin y)\,\mathrm{d}\sigma=0.
      \]

    \item \textbf{Solution}. B.

      正弦级数表示的是 \([0,\pi]\) 上函数的奇延拓. 因而
      \[
        S\left(-\frac{\pi}{2}\right)=-f\left(\frac{\pi}{2}\right)
        =-\left(1-\frac{\pi}{2}\right)
        =-\frac{\pi}{2}+1.
      \]

    \item \textbf{Solution}. C.

      由对称性
      \[
        \oint_\Gamma y\,\mathrm{d}s=\oint_\Gamma z\,\mathrm{d}s=0,
      \]
      故
      \[
        \oint_\Gamma(y+2z)\,\mathrm{d}s=0.
      \]

    \item \textbf{Solution}. A.

      若 \(\sum u_n\) 收敛，则必有 \(u_n\to 0\)，从而充分大时 \(|u_n|\le 1\)，于是
      \[
        u_n^2\le |u_n|.
      \]
      由比较判别法，\(\sum u_n^2\) 收敛.
  \end{enumerate}

  \section{填空题(每小题 4 分，共 16 分)}
  \begin{enumerate}
    \setcounter{enumi}{6}
    \item \textbf{Solution}. $2\pi$.

      因为曲线 \(\Gamma\) 位于平面 \(x+y+z=0\) 上，且关于三个坐标完全对称，所以
      \[
        \oint_\Gamma x^2\,\mathrm{d}s=\oint_\Gamma y^2\,\mathrm{d}s=\oint_\Gamma z^2\,\mathrm{d}s.
      \]
      又
      \[
        \oint_\Gamma(x^2+y^2+z^2)\,\mathrm{d}s=\oint_\Gamma 1\,\mathrm{d}s=2\pi.
      \]
      从而
      \[
        \oint_\Gamma(x^2+2y^2+3z)\,\mathrm{d}s=2\pi.
      \]

    \item \textbf{Solution}. $x+3y-z=0$.

      由
      \[
        z=x+3y+\ln(1+x^2+y^2)
      \]
      知
      \[
        z_x(0,0)=1,\qquad z_y(0,0)=3.
      \]
      在点 \((0,0,0)\) 处的切平面方程为
      \[
        z=x+3y.
      \]

    \item \textbf{Solution}. $\dfrac19$.

      \[
        u=\ln\sqrt{x^2+y^2+z^2}=\frac12\ln(x^2+y^2+z^2).
      \]
      在三维空间中
      \[
        \Delta u=\operatorname{div}(\operatorname{grad}u)=\frac{1}{x^2+y^2+z^2}.
      \]
      代入 \((1,-1,2)\) 得
      \[
        \Delta u=\frac{1}{1+1+4}=\frac19.
      \]

    \item \textbf{Solution}. $\dfrac{\pi}{2}$.

      傅里叶级数的和函数是原函数的 \(2\pi\) 周期延拓在间断点取左右极限平均. 因为
      \[
        -3\pi\equiv -\pi\equiv \pi \pmod{2\pi},
      \]
      故
      \[
        S(-3\pi)=\frac{f(\pi^-)+f(-\pi^+)}{2}
        =\frac{2\pi+(-\pi)}{2}
        =\frac{\pi}{2}.
      \]
  \end{enumerate}

  \section{基本计算题(每小题 7 分，共 42 分)}
  \begin{enumerate}
    \setcounter{enumi}{10}
    \item \textbf{Solution}. 与直线
      \[
        L_1:\ \frac{x-1}{2}=\frac{y-2}{1}=\frac{z-2}{-1}
      \]
      垂直的平面法向量可取 \((2,1,-1)\). 过点 \(P(2,1,3)\) 的平面方程为
      \[
        2(x-2)+(y-1)-(z-3)=0,
      \]
      即
      \[
        2x+y-z-2=0.
      \]
      将
      \[
        x=2t+1,\quad y=t+2,\quad z=-t+2
      \]
      代入该平面方程，可得 \(t=0\)，故交点为
      \[
        (1,2,2).
      \]

    \item \textbf{Solution}. 先求
      \[
        z_y=e^{-x}\cos\frac{y}{x}\cdot \frac{1}{x}=\frac{e^{-x}}{x}\cos\frac{y}{x}.
      \]
      所以
      \[
        z_y(2,\pi)=\frac{e^{-2}}{2}\cos\frac{\pi}{2}=0.
      \]
      再对 \(x\) 求偏导，可得
      \[
        z_{yx}=e^{-x}\left(\frac{y}{x^2}\sin\frac{y}{x}-\frac{x+1}{x^2}\cos\frac{y}{x}\right)\frac{1}{x}.
      \]
      在 \((2,\pi)\) 处，\(\sin\dfrac{\pi}{2}=1,\ \cos\dfrac{\pi}{2}=0\)，故
      \[
        z_{yx}(2,\pi)=\frac{\pi}{8e^2}.
      \]

    \item \textbf{Solution}. 曲面面积为
      \[
        S=\iint_{x^2+y^2\le 1}\sqrt{1+z_x^2+z_y^2}\,\mathrm{d}x\,\mathrm{d}y.
      \]
      由于 \(z=xy\)，故
      \[
        z_x=y,\qquad z_y=x.
      \]
      所以
      \[
        S=\iint_{x^2+y^2\le 1}\sqrt{1+x^2+y^2}\,\mathrm{d}x\,\mathrm{d}y.
      \]
      用极坐标计算：
      \[
        S=2\pi\int_0^1 r\sqrt{1+r^2}\,\mathrm{d}r
        =\frac{2\pi}{3}(2\sqrt2-1).
      \]

    \item \textbf{Solution}. 绕 \(z\) 轴旋转后得到椭球体
      \[
        x^2+y^2+\frac{z^2}{2}\le 1.
      \]
      对固定 \(z\)，截面是圆
      \[
        x^2+y^2\le 1-\frac{z^2}{2},
      \]
      面积为
      \[
        \pi\left(1-\frac{z^2}{2}\right).
      \]
      因而
      \[
        I=\int_{-\sqrt2}^{\sqrt2} z^2\pi\left(1-\frac{z^2}{2}\right)\,\mathrm{d}z
        =\frac{8\sqrt2\pi}{15}.
      \]

    \item \textbf{Solution}. 将上半球与底圆盘补成闭曲面，应用高斯公式. 由对称性，含 \(y\) 的奇函数部分积分为零，最后只需处理对称项，计算得
      \[
        I=\frac{2\pi R^2}{3}.
      \]

    \item \textbf{Solution}. 用比值法可知收敛半径为 \(\infty\). 又
      \[
        S(x)=\sum_{n=1}^{\infty}\frac{2n+1}{n!}x^{2n}
        =2\sum_{n=1}^{\infty}\frac{n}{n!}x^{2n}+\sum_{n=1}^{\infty}\frac{x^{2n}}{n!}.
      \]
      其中
      \[
        \sum_{n=1}^{\infty}\frac{n}{n!}x^{2n}=x^2e^{x^2},\qquad
        \sum_{n=1}^{\infty}\frac{x^{2n}}{n!}=e^{x^2}-1.
      \]
      所以
      \[
        S(x)=(2x^2+1)e^{x^2}-1.
      \]
  \end{enumerate}

  \section{综合题(每小题 7 分，共 14 分)}
  \begin{enumerate}
    \setcounter{enumi}{16}
    \item \textbf{Solution}. 记
      \[
        P(x,y)=(-xe^{-x}+f''(x))y,\qquad Q(x,y)=f'(x).
      \]
      由路径无关得
      \[
        P_y=Q_x,
      \]
      即
      \[
        -xe^{-x}+f''(x)=f'(x).
      \]
      于是
      \[
        f''(x)-f'(x)=xe^{-x}.
      \]
      解该微分方程并利用
      \[
        f(0)=4,\qquad f'(0)=3
      \]
      可得
      \[
        f(x)=4e^x+\frac12(x^2-x)e^x.
      \]

    \item \textbf{Solution}. 在区域内部，
      \[
        f_x=2x,\qquad f_y=-2y,
      \]
      唯一驻点是 \((0,0)\)，其函数值为 \(2\).

      在边界
      \[
        4x^2+y^2=4
      \]
      上有 \(y^2=4-4x^2\). 代入原函数得
      \[
        f=x^2-(4-4x^2)+2=5x^2-2,\qquad -1\le x\le 1.
      \]
      所以边界上最大值为 \(3\)，最小值为 \(-2\). 与内部值比较知
      \[
        \max f=3,\qquad \min f=-2.
      \]
  \end{enumerate}

  \section{证明题(每小题 5 分，共 10 分)}
  \begin{enumerate}
    \setcounter{enumi}{18}
    \item \textbf{Proof}. 交换积分次序，有
      \[
        \int_0^1\,\mathrm{d}x\int_x^1 f(x)f(y)\,\mathrm{d}y
        =\int_0^1\,\mathrm{d}y\int_0^y f(x)f(y)\,\mathrm{d}x.
      \]
      将两式相加：
      \[
        2\int_0^1\,\mathrm{d}x\int_x^1 f(x)f(y)\,\mathrm{d}y
        =\int_0^1\int_0^1 f(x)f(y)\,\mathrm{d}x\,\mathrm{d}y.
      \]
      右边可分离为
      \[
        \left(\int_0^1 f(x)\,\mathrm{d}x\right)\left(\int_0^1 f(y)\,\mathrm{d}y\right)=A^2.
      \]
      因而
      \[
        \int_0^1\,\mathrm{d}x\int_x^1 f(x)f(y)\,\mathrm{d}y=\frac{A^2}{2}.
      \]

    \item \textbf{Proof}. 由递推式
      \[
        2a_{n+1}=a_n+\sqrt{a_n^2+u_n}
      \]
      可知 \(\sqrt{a_n^2+u_n}>a_n\)，因此
      \[
        a_{n+1}>a_n,
      \]
      即 \(\{a_n\}\) 单调递增. 题设又给出它有上界，所以 \(\{a_n\}\) 收敛.

      再由递推式整理得
      \[
        u_n=4a_{n+1}(a_{n+1}-a_n).
      \]
      因为 \(a_n>0\) 且单调递增，所以
      \[
        0<u_n\le 4(a_{n+1}^2-a_n^2).
      \]
      而
      \[
        \sum_{n=1}^{\infty}(a_{n+1}^2-a_n^2)
      \]
      为望远镜级数，并且由于 \(\{a_n\}\) 收敛而有有限和. 由比较判别法，
      \[
        \sum_{n=1}^{\infty}u_n
      \]
      收敛.
  \end{enumerate}
\end{document}






